\documentclass[12pt, twoside]{report}
\usepackage[utf8]{inputenc}
\usepackage{graphicx}
\graphicspath{ {images/} }
\usepackage{caption}
\usepackage{subcaption}
\usepackage{xcolor}
\usepackage{gensymb}
\usepackage[a4paper, width=150mm, top=25mm, bottom=25mm]{geometry}
\usepackage[backend=biber]{biblatex}
\addbibresource{references.bib}



\begin{document}

\begin{titlepage}
	\begin{center}
	
	\vspace*{0.5cm}
	
	\Huge
	\textbf{Toward Scalable and Robust Agent Based Methods.}
	
	\vspace{0.5cm}
	\Large
	An Analysis of Active Inference and Reinforcement Learning Paradigms in Large, Partially Observable and Non-Statioanry Environments.
	
	\vspace{1.5cm}

	\textbf{Fraser Paterson}

	\vspace{1.5cm}

	A Thesis proposal, pursuant to the requirements of the\\ 
	Degree: Bachelor of Science (Honours).  
	
	\vspace{2.0cm}

	\includegraphics[width=0.4\textwidth]{UWA_Logo.png}
	
	\vspace{2.0cm}	
	
	\Large
	Supervisor: Dr Tim French\\ 
	Department of Computer Science and Software Engineering\\
	The University of Western Australia\\
	12 March 2023
	
	\end{center}
\end{titlepage}


\tableofcontents

\cleardoublepage



\section{Introduction}

\subsection{Motivation and Background}
Many real-world problems are characterised by high degrees of noise, ill-definedness and uncertainty. This uncertainty assumes myriad forms, whether in the clarity of the observations one can solicit from the system of interest (signal detection), or in the confidence of an inference as to the system-parameter values, that best account for the solicited observation (state-estimation). These tasks are only further complicated by a very common constraint on any candidate solution technique: partial-observability, which is an obvious source of uncertainty. Indeed, partial-observability is overwhelmingly characteristic of many ``difficult'' real-world systems. 

It is the perogative of the Engineer to make simplifying assumptions about the kinds and sources of uncertainty which beset the task of modeling, controling and/or predicting the behaviour of a system. A ``good'' simplification (model) of a complicated system, affords some relevant trade-off between the computational resources necessary to predict and control the entities of interest, as against the material and temporal constraints necessitated by the nature of the problem. This trade-off is inherent to cognitive agency itself; it is to be in the ``Finitary Predicament'', see: \textcite{Cherniak}. For example, a ``pedestrian-detection'' algorithm on a self-driving car, plausibly needs to put a premium on speed since temporal constraints are quite salient in this case. An algrithm for optimising the trajectory of a spacecraft mission 5 years from now faces the same trade-off, but under near-inverse constraints. In this case, it is plausable to suppose that there is not such a high premium on time. In any case, a perennial problem remains, that of obtaining accurate \textit{enough} observations and performing an optimal \textit{enough} action - at the optimal \textit{enough} location and time - for resolving the maximal ammount of uncertainty about the dynamics one wishes to predict and/or control, while also minimising the costs associated with planning and performing these actions.  

The faculty by which complex biological organisms achieve these ends is by means of their intelligence. A relatively non-controvercial definition of \textit{Intelligence} is ``general problem-solving ability''. The ongoing project of defining, detecting and creating instances of intelligence is one that takes a startling variety of forms. Psychologists speak of ``intelligent behaviour'', whereas to the Neuroscientist, intelligence is studied in terms of ``patterns of neuronal activation''. Yet again, the computer scientist traditionally accounts for intelligence as a kind of ``information processing''. Insofar as these disciplines might share a common project, it may plausibly be supposed as the attempt to understand the nature and function of intelligence. Several philosophical approaches to this project abound - as attested by the variety of disciplines it spans. Again, the Psychologist studies \textit{behaviour}, the neuroscientist studies \textit{the brain} but the computer scientist has an altogether distinct modality of investigation. This ``investigative modality'' consists in the attempt to build, create and/or instantiate an \textit{instance} of intelligence, hence ``Artificial'' Intelligence. The hope is that if we are succesful enough at reproducing the effects of natural intelligence, that this will count as a good reason to suppose that the nature and function of our artificial intelligence, is constitutive of the nature and function of ``Natural Intelligence''. Thus the project of Artificial Intelligence consists in the attempt to analyse, formalise and mechanise Natural Intelligence. 

An early and perhaps paradigm case of an attempt to so analyse, formalise and mechanise intelligence, was Newell and Simon's General Problem Solver (GPS): \textcite{Newell-Simon}.
While the GPS could reliably reproduce known proofs for various therorems in Mathematics and Logic, it failed to provide solutions to less well-deffined problems. Newell and Simon had assumed that all problems were essentially the same and thus that the formulation of a problem was a relatively trivial task. This turns out to have been false, all problems are not essentialy the same, indeed most of the ``difficulty'' associated with the truly vexing, non-contrived sorts of ``real world'' problems of practical interest, arises from the question as to how we may best represent the problem. see \textcite{AFTMT-Problem-Formulation} for more on the GPS, its limitations and its legacy to the project of Artificial Inteligence. The insight that the GPS formalisation afforded was crucial. Even if we assume that we have a problem forulation,the vast majority of all search spaces are exponential in nature. It is literally impossible for the sorts of epistemically-bounded creatures that we are - and that we seek to create - to check the whole space of solutions, or even to check most of it. Thus we must somehow zero-in on the \textit{relevant} information, and intelligently \textit{ignore} the vast majority of the search space. see: \textcite{RR-PP} and \textcite{RR-Emerging} for an ellaboration. Thus we see that any attempt to create genuine \textit{instances} of intelligence must be furnished with some means to bias their attention in the search space of possible solutions. 

A common means by which the Engineer has attempted this task is by constructing model-based artificial agents. A simplified model of the system of interest affords an agent the ability to bias its attention to those parts of the system that are likely to be relevant to the task at hand. As an example, a model affords the ability to evaluate counter-factual claims about the environment: ``what would it be like if I did X at Y?'' and thus afford the ability for an agent to \textit{plan} actions in its environment.      

Model based artificial agents have enjoyed great sucess in recent years: \textcite{Silver2017} and \textcite{dream}, owing to the advantages outlined above (among other things). Indeed, greater sample-efficency over their model-free cousins is a direct consequence of the model's ability to bias attention toward the most relevant trajectories through the search space. Model based methods are no cure-all, they have trade-offs of their own, and in that respect, they are like every other method.

At least two canonical problems exist for any model-based method. The first concerns the degree to which the model can continue to provide robust predictions, in the face of various sources of uncertainty. This uncertainty may be in the form of noisy state-observations, partial observability or the existence of non-stationary environmental parameters. The second problem concerns the degree to which the model can be used to provide apt predictions in higher-dimensional environments, that is, how best to ``scale up'' the model to a larger problem. These issues are related, often an attempt to ``scale up'' a working model will increase the model's complexity and this very often negatively impinges upon the model's robustness.    

Indeed an exact Bayesian approaches to optimal inference is almost always intractible, due to the necessity for marginalisation over all states (in what is very often, an exponential search space). This is computationally intractible for any but the simplest search spaces. Most real-world environments are continuous, high-dimensional and non-stationary. These issues are instantiations of two broard ``problem classes'', though there is overlap between the two. These two classes are modle robustenss and scalibility. In other words, how well can the model operate in the presence of uncertainty? and how well can the approach generalise to higher-dimensional spaces?

Before ellaborating upon the central research objectives of this proposed thesis, it will be necessary to provide some background on Active Inference methods, and to justify their interest as plausible agent-based methods.


\subsection{Active Inference: An Overview}
Active Inference (AIF) is a corollary of the Free Energy Principle (FEP), as it pertains to the imperatives attendent upon embodied, cognitive agents. A few words about the FEP shall suffice. The FEP is itslef born out of the interface between Statistical Mechanics, Infomation Theory and Cognitive Science and can trace its roots back to the work of Gibbs, Hamilton and Helmholtz (among many others). Largely popularised by the work of Karl Friston and others, the FEP is a plausible, unifying account of brain function in which the brain is supposed as engaging in a scale-invariant process of Variational Free Energy minimisation - over sensory data - so as to maximise its own model-evidence (thereby resisting a thermodynamic tendency to dissolution and ultimately, death). See: \textcite{FEP-Rough-Guide-Brain}, \textcite{Action-Behaviour-FE} and \textcite{FEP-Math-Review} for details. 

The Variational Free Energy: $\mathcal{F}$ is a functional of beliefs over - potentially uncertain - sensory observations. $\mathcal{F}$ is provably a lower bound, on an information-theoretic quantity called ``surprisal''. The mathematical expression for surprisal is: $-\log{p(o)}$ where $o \in O$ is an observation. Hence we have $$\mathcal{F} \geq -\log{p(o)}$$ See: \textcite{FEP-Rough-Guide-Brain}, \textcite{Action-Behaviour-FE} and \textcite{FEP-Math-Review} for a full derivation of this result, and for more on the FEP's relation to Active Inference.    

Surpirsal quantifies the ``atypicallity'' or ``unexpectedness'' of an observation, with respect to the predictions afforded by the generative model. Since the Free Energy is a lower bound on sensory suprisal, if we can merely minimise this VFE, we shall have implemented a means of approximate Bayesian infernce. This method of approximating the true posterior, by means of a variational distribution is well known. Indeed it is called ``Variational Inference'' - predictably enough. See \textcite{VI} for more. 

Now while the above relations afford an agent the means to approximate surprisal, the mere ability to evaluate the free energy does nothing to minimise surprisal. The agent can only minimise its surprisal vicariously. To this end, there are the following choices. The agent can either change the structure of its generative model so as to better conform to its present observations (a process of perception) or it can act on its environment, so as to change its observations, such that they better accord with its generative model (a process of action). Hence ``Active'' Inference. In a nutshell, if there is a discrepancy between your model of the world and your observations about the world, you can either try and change your model, or you can try and change the world.  

This contrasts the Reinforcement Learning paradigm, where an agent learns a state-to-action mapping that maximise a scalar reward-signal over some potentially infinite time-horizon.
Indeed an AIF differes in another funcamental respect to Reinforcement Learning. In AIF, there is a natural imperative that drives the agent to perform actions which are maximally ``uncertainty reducing''. This is a result of the mathematical form of the VFE. By measureing the seperation between the true posterior over environmental states: $P(s|o)$ and the approximate, variational posterior: $Q(s)$ as the KL Divergence between the two, we can derive an expression for the Variational Free Energy, $\mathcal{F}$. $$D_{KL}(Q(s) || P(s|o)) = \int Q(s)\log(\frac{Q(s)}{P(s | o)})ds$$ This can be re-expressed as $$D_{KL}(Q(s) || P(s|o)) = \mathcal{F} + \log{P(o)}$$ where $$\mathcal{F} = \int Q(s)\log(\frac{Q(s)}{P(s, o)})ds = D_{KL}(Q(s) || P(o, s))$$

Hence we have $$\mathcal{F}  = D_{KL}(Q(s) || P(s|o)) - \log{P(o)}$$ The first term acts as an epistemic imperative, driving the agent to perform those actions which maximally reduce the discrepancy between its approximate beliefs and the true dynamics. The latter term affords a pragmatic imperative, in that the agent is driven to perform actions in such a way so as to minimise the probability of encountering ``surprising'' states. This surpirsal has a direct analog to reward in the Reinfircement Learning context but instead of being a scalar signal, it is simply a modality specific observation over which the agent has some preference.  


\section{Research Objectives}
Real-world artificial agents face two serious issues. The first is a high degree of uncertinty, whether this be a consequence of noisy observations, partial observability or non-stationary environmental dynamics. Uncertainties of this kind can reduce an otherwise effective agent to an utter failure. Thus, is the issue of affording \textbf{robust} agents. The second issue concerns the nature of difficult, real-world environments, in that they are very often high-dimensional and/or continuous. Simple exploration/planning algorithms can quickly be rendered ineffective when situated in a higher-dimesnioanl instantiation. This issue is known as \textbf{scaling}.

Reinforcement Learning has enjoyed a great deal of success in the attempt to scale to higher dimensional, continuous, noisy environments - examples. While there is still much to be done on this front, Active Inference has thus far been almost entierly limited to small, discrete environments, see: \textcite{AIF-D}, \textcite{AIF-Cur-Insight}. Given that Active Inference represents a potentially unifying paradigm - owing to its generality - and that it has no dependence on any ad-hoc scalar reward signal, it is plausible to suppose that Active Infernce might enjoy several theoretical advantages over more ``traditional'' methods in Reinforcement Learning and Optimal Control, see: \textcite{RL-or-AIF}, \textcite{Friston2012} and \textcite{Optim-Motor}.   

Hence, in the course of this proposed thesis, we shall implement various agents - AIF and RL agents - in both the the fully observable and partially observable cases, with and without various sources of uncertainty and in both benchmark ``low-diensional'' environments and  non-trivial ``higher-dimensional'' environments. The aim will be to asses the relative advantages and disadvantages of both ``approach families'' as these pertain to their ability to effectively deal with the various kinds of uncertrainty mentioned above, in addition to their ability to scale up into higher-dimensional environments.     


\subsection{Central Research Questions}
The principle aim of this research thesis is twofold. The first aim is to investigate the robustness of AIF methods in noisy, uncertain or partially observable environments, relative to Reinforcement Leraning baselines. The second aim concerns the potental for ``scaling up'' AIF methods to continuous and/or higher-dimensional state-spaces. 

Thus are the central questions asked in this Thesis:

\begin{itemize}
\item Are AIF agents more robust to noisy observations and non-stationarity than a comparable RL baseline?
\item Can AIF be scaled up to higher dimensional Partially Observable envronments? 
\end{itemize} 



\section{Previous Work}
Work of this kind has already appeared in the literature, though it is still very much in its infancy. \textcite{Markovi-2021} implemented an Active Inference agent for the multi-armed bandit problem, in the stationary and non-stationary case. The AIF agent did not perform as well as a state-of-the-art bayesian UCB algorithm, in the stationary case. However in the non-stationary case, the AIF agent outperformed the UCB agent. While this implementaion was over a small, discrete space of environmental states, the results plausibly suggest that AIF would be an effective means of optimal inference and control in a higher-dimensional or continuous problem.  

An approach that has enjoyed some success as of late involves the implementaion of Free Energy minimisation as a process of message-passing on a Forney-style factor graph. See: \textcite{Sim-AIF-Message}, \textcite{Cox-2019}, \textcite{Reactive-MP} and \textcite{Deep-Temp-AIF}. In this framework, the agent's generative model is factorised in such a was as to be a Forney or ``Normal'' factor graph. Free Energy minimisation is then cast as a process of message passing over this factor graph. Various message passing algorithms exist, such as Belief Propagation and Variational Message Passing. This message passing scheme greatly reduces the number of terms over which it is necessary to sum, when computing the approximate marginal and posterior distributions thereby affording tractible Active Inference algorithms in relatively high-dimensional settings. The natural inclusion of a generative model and the built-in epistemic imperatives - toward the aim of uncertainty reduction - in AIF, make it highly plausable that this method will be better able to deal with non-stationary environments, dynamic constraint changes, noise and other such sources of uncertainty. Indeed, just as much has been shown in \textcite{Bandits}, where AIF performd better than a strong Bayesian UCB algorithm in a non-stationary multi-armed bandit problem.

Yet other approaches have atempted to leverage the ability for deep neural networks to parameterise the distributions of interest. See: \textcite{Deep-AIF-Ueltzh-2018}, \textcite{Deep-Var-Policy-Grad} and \textcite{DEEP-AIF-POMDPs}. Of particular intetest with this aproach are \textcite{Scaling-AIF} and \textcite{Contrastive-AIF}. The former makes use of amortized inference, in the form of neural network function approximators to parameterize the relevant distributions. Free Energy minimisation is then performed with repect to the function approximators. The use of amortized inference affords several advantages. For example, the number of parameters remains constant with respect to the size of the data and inference can be achieved via a single forward pass through the network. The resulting algorithm was able to explore a much greater proportion of the state space in a simple MountainCar environment, as opposed to two Reinforcement Learning, baseline agents. In addition, the agent was able to learn to control the continuous inverted pendulum task with a far greater sample efficency than the baseline agents. Although the approach offered in \textcite{Scaling-AIF} is promising, it was restricted in every case to fully observable environments. 

Lastly, the approach of \textcite{Contrastive-AIF}, implemented a contrastive objectove for their Active Inference agent, which signifigantly reduced the computational burden of learning the generative model and planning future actions. This method performed signifigantly better than the usual, liklihood-based ``reconstructive'' means of implementing AIF and it was also computationally cheaper to train. Importantly, this method offered a unique way to afford increased model-robustness in the face of environmental distractors. 


\section{Methods and Objectives}
The proposed structure of the investigation shall obey the follwoing three-part itinerary. 

\subsection{Benchmark AIF and RL Agents}
Initially, our analysis will be limited to the fully observable case. The environments of interest will be the well-known MountainCar and CartPole baseline environments. The aim will be to assess the performance - robustness, learning rate and sample eficency - of the AIF and Reinforcement Learning agents across two dimensions. These two dimensions are as follows:
\begin{enumerate}
\item Whether the AIF agent is implemented in OpenAI gym or RxInfer.jl (testing performance diferences in the factor-graphical method)
\item Whether the environment is ``uncertain'' (characterised by either noisy observations or non-stationarity)
\end{enumerate}

The AIF agent will be instantiated by means of a deep generative model, and the RL baseline agent will be DQN. The stochastic version of each environment will either furnish the agent with observations that have been subject to additive Gaussian white noise, or will be characterised by a degree of non-stationarity. An example of a potential kind of non-statioarity in the CartPole environment might include randomly changing the length of the pole at iteration n, for all remaining episodes - for instance.  

Research questions at this level shall include:
\begin{itemize}
\item What is the sample efficency of the two approach families?
\item What is the complexity of the two approach families?
\item Does the Factorised generative model appreciably effect the performance of the AIF agent as against the baseline Rl agent?
\item Is one approach family - AIF or RL - ultimately more robust to noise or non-stationarity?
\end{itemize}  
   
\subsection{Deep AIF - The Partially Observable Case}
This portion of the study will take up the investigation offered in \textcite{Scaling-AIF} whereby we shall attempt to implement their Deep AIF agent, only now in the partially observable case. 
We shall attempt to implement a Deep AIF Pong plaing agent, from the Atari library in OpenAI gym. Pong is technically a partialy observable environment, since a single frame does not afford any information about the velocity of the ball. Again, we shall assess the performance of the Deep AIF as agianst a DQN baseline, in the stationary and non-stationary POMDP.  

Research questions at this latter level will echo those of the latter:
\begin{itemize}
\item What is the sample efficency of the two approach families? has it appreciably changed?
\item What is the complexity of the two approach families?
\item Is one approach family - AIF or RL - ultimately more robust to noise or non-stationarity?
\end{itemize}
  
\subsection{Deep Contrastive AIF With Partial Observability}
Time permitting, and assuming that the latter two tasks are achieved with dome degree of satisfacation, the analysis will be extended finally to a much higher-dimensional POMDP such as Asteroids. The Deep AIF agent will be extended to use the contrastive learning approach developed in \textcite{Contrastive-AIF}. This approach signfigantly reduces the computational complexity of the model and hence, appears to be a promising avenue of investigation in the extension of Deep AIF methods to high-dimensional, partially observable models.   

\section{Software and Hardware Requirements}
Most of the simulations will be instantiated in Python and associated libraries  - such as OpenAI gym.
RxInfer.jl will be used to implement the factor-graph based AIF agents in the first portion of the study.

The Deep Versions of each agent can be more effectively trained by means of an Nvidia 3060 GPU, available to the author.  


\printbibliography

\end{document}
